%% =====================================================================
%%  T-13-03  The Veil
%%  -------------------------------------------------------------------
%%  One idea per file. Core is mandatory. Slots in canonical order.
%%  Max 700 words. No layout commands. No filler. New source material
%%  for this entry goes in sources/B3/T-13-03.tex -- never by
%%  rewriting this file (docs/RULES.md R0.1).
%%  =====================================================================
%% @kind: topic
%% @id: T-13-03
%% @title: The Veil
%% @subtitle: Conceal intentions, say less --- the unreadable opponent cannot be prepared for
%% @chapter: c13
%% @order: 30
%% @status: stable
%% @sources: B3
%% @updated: 2026-09-19

\topic{T-13-03}{The Veil}{Conceal intentions, say less --- the unreadable opponent cannot be prepared for}

\begin{core}
Two veils, one principle. Conceal the purpose behind your actions: people off balance and in the dark cannot prepare a defence, and the moment your intention is named --- even by one careless word --- every move you made is re-read in its light. And say less than necessary: the more you say, the more common you appear and the more material you hand over. The unreadable opponent is the un-preparable-for opponent.
\end{core}
\begin{mechanism}
B3 grounds the first veil in human nature's first instinct: we trust appearances, so the sincere-seeming surface is accepted while the plan runs beneath it --- false sincerity, decoy objects of desire, red herrings across the scent, the bland exterior that hides what a poker face would advertise. The mechanism the veil exploits is single focus: people can attend to one thing at a time, and attention directed at the staged object is attention not pointed at the real one. The second veil is verbal economy. Overtalking is information transfer plus status leak: every extra sentence narrows the target's uncertainty about you, and banality dressed in volume reads as neediness. The Oracle works in the opposite direction --- few enigmatic words, weighted by the listener's own imagination. Silence also forces the counterpart to fill it: nervous talk under your quiet is the cheapest intelligence collection there is (T-07-01). The two veils share a failure mode, which the source states as a reversal: no screen conceals anything once you carry a reputation for deception --- the veil is a spend-it-slowly asset.
\end{mechanism}
\begin{conditions}
Strongest in adversarial or negotiational settings with strangers and rivals, and in any long game whose value dies with early disclosure. It weakens among intimates, where concealment is corrosive rather than protective, and in institutions built on disclosure --- sworn testimony, audited accounts --- where the veil's discovery is the catastrophe. Against professional readers, the veil itself is a tell: too much blandness is a pattern (T-23-02).
\end{conditions}
\begin{application}
  \item Separate the plan from the presentation. Decide what surface the room should see, and rehearse the surface.
  \item Give decoys somewhere to point: a visible want that shares the budget of attention with the real one.
  \item Cut your speaking part by half, then cut the explanations. Say the conclusion; withhold the derivation.
  \item Use silence as a question. The pause after their statement collects what their defence prep cannot survive.
  \item Spend the veil slowly. One revealed intention buys years of credibility for the next concealment --- budget it.
\end{application}
\begin{feedback}
  \item Working: rivals prepare for the decoy. Resources spent defending the wrong object are a gift.
  \item Working: people leave conversations with you unsure, and interest rises with the uncertainty.
  \item Failing: one confessed motive re-prices every past move --- the veil is seen and everything behind it re-audited.
  \item Failing: your brevity reads as coldness in a relationship that runs on disclosure. Wrong room for the veil.
\end{feedback}
\begin{failure}
The veil's occupational hazard is the operator's own map: conceal so long and you start hiding from yourself --- intentions you have never had to state are intentions you cannot inspect. The silence weapon also has a domestic price: relationships run on derivations, not conclusions, and the sphinx at home is just absent. And veils are contagious: run one on a partner and they will run one back --- the game then has no information in it at all.
\end{failure}
\begin{countermeasures}
  \item When someone is unfailingly bland and agreeable, look for the decoy: what are you being shown, and what is being bought with your attention?
  \item Demand derivations, not conclusions, from anyone selling certainty. Veils live on conclusions.
  \item Track disclosure asymmetry in your dealings: if they know your intentions and you know nothing of theirs, you are not negotiating, you are being harvested (T-07-01).
  \item Fill silence on purpose, never by accident. Decide before meetings what you will say into a pause --- and what you will not.
\end{countermeasures}

\sources{B3}
\seesources
\seealso{T-13-02, T-07-01, T-16-03}
